\section*{第 \ref{ch:g10:drainage-basins} 章——流域、河流与洪水}

\begin{solution}{exo:g10:drainage-basins:1}
河流某一点的\omterm{def:g10:drainage-basins:drainage-basin}{流域}，是指这样一片陆地的全部范围，其径流
流入该点上游的河段并经过该点。\omterm{def:g10:drainage-basins:drainage-basin}{分水岭}
是沿山脊延伸、分隔两个相邻\omterm{def:g10:drainage-basins:drainage-basin}{流域}的界线。
\omterm{def:g10:drainage-basins:drainage-basin}{分水岭}两侧相距一米落下的两滴水，沿相反的
坡面流入两个不同的\omterm{def:g10:drainage-basins:drainage-basin}{流域}：它们的河流可以流入不同的海。
\end{solution}

\begin{solution}{exo:g10:drainage-basins:2}
\omterm{def:g10:drainage-basins:discharge}{流量}是指河流某一剖面上每秒流过的水体体积，
单位为 \unit{m^3/s}。\qty{1}{m^3/s} 即
每秒 \qty{1000}{L}；一天共 $1000 \times 86\,400 = 86.4$ 百万升。
\end{solution}

\begin{solution}{exo:g10:drainage-basins:3}
洪水来临前，河流\omterm{def:g10:drainage-basins:discharge}{流量}约为 \qty{330}{m^3/s}（介于约
\num{240} 与 \qty{390}{m^3/s} 之间）。洪峰\omterm{def:g10:drainage-basins:discharge}{流量}为 \qty{9231}{m^3/s}，出现在
1 月 21 日。涨水始于 1 月 18 日（\qty{362}{m^3/s}），历时三
天。
\end{solution}

\begin{solution}{exo:g10:drainage-basins:4}
$1/50 = 0.02$，即 \qty{2}{\percent}。不对：各年彼此独立，
无论去年如何，今年的概率仍为 \qty{2}{\percent}。
\end{solution}

\begin{solution}{exo:g10:drainage-basins:5}
1 立方英尺为 $0.3048^3 \approx \qty{0.02832}{m^3}$，因此
$593\,000 \times 0.02832 \approx \qty{16800}{m^3/s}$。亚马孙河的\omterm{def:g10:drainage-basins:discharge}{流量}
约为其 $206\,000/16\,800 \approx 12$ 倍。
\end{solution}

\begin{solution}{exo:g10:drainage-basins:6}
条带面积：$1.5 \times 0.3 = 0.45$、$0.90$、$1.05$ 和 \qty{0.60}{m^2}，
总面积 $A = \qty{3.0}{m^2}$。各条带\omterm{def:g10:drainage-basins:discharge}{流量}：$0.2 \times 0.45 = 0.09$、$0.5 \times 0.90
= 0.45$、$0.6 \times 1.05 = 0.63$、$0.3 \times 0.60 = 0.18$，总\omterm{def:g10:drainage-basins:discharge}{流量} $Q =
\qty{1.35}{m^3/s}$。平均流速为 $1.35/3.0 = \qty{0.45}{m/s}$。
\end{solution}

\begin{solution}{exo:g10:drainage-basins:7}
$66 \times 10^6 \times 0.020 = \qty{1.32e6}{m^3}$。两天的时长
为 $2 \times 86\,400 = \qty{172800}{s}$，$1.32 \times 10^6/172\,800 \approx
\qty{7.6}{m^3/s}$，即平均值。
\end{solution}

\begin{solution}{exo:g10:drainage-basins:8}
由大到小排列：120、88、65、59、51、47、42、38、35。取 $n = 9$，$T = 10/m$，则
最大值（120）的 $T = 10$ 年，第二大（88）的 5 年，第三大（65）的
约 3.3 年。中位数 51 排第 5，$T = 2$ 年：
平均每两年即有一年达到或超过该值。
\end{solution}

\begin{solution}{exo:g10:drainage-basins:9}
城镇的\omterm{def:g10:drainage-basins:hydrograph}{流量过程线}涨得更早、更高，回落也很快：在
陡坡、屋顶和道路上，入渗的水很少，径流迅速
且集中地到达溪流，其\omterm{def:g10:drainage-basins:hydrograph}{滞后时间}很短。林区\omterm{def:g10:drainage-basins:drainage-basin}{流域}的
\omterm{def:g10:drainage-basins:hydrograph}{流量过程线}涨得较晚，洪峰较低、较宽：土壤和森林
吸收了大部分降雨并缓慢释放，水流在平缓的坡面上
流动缓慢。风暴过后，其\omterm{def:g10:drainage-basins:hydrograph}{基流}仍保持较高。
\end{solution}

\begin{solution}{exo:g10:drainage-basins:10}
$9231 \times 86\,400 \approx \qty{7.98e8}{m^3}$，
若均匀铺展在 $29\,900 \times 10^6 = \qty{2.99e10}{m^2}$ 的面积上，约
为 \qty{0.027}{m}，即 \qty{27}{mm}。\omterm{def:g10:drainage-basins:hydrograph}{基流}给出 $330 \times 86\,400/2.99 \times 10^{10}
\approx \qty{0.001}{m}$，约合每天 \qty{1}{mm}：洪峰当天的水量约为
其 28 倍。
\end{solution}

\begin{solution}{exo:g10:drainage-basins:11}
$P = 1 - 0.9^{10} \approx 1 - 0.349 = 0.651$，约为 \qty{65}{\percent}。
完全不发生的概率为 $0.9^{10} \approx 0.35$，约
\qty{35}{\percent}。
\end{solution}

\begin{solution}{exo:g10:drainage-basins:12}
该\omterm{def:g10:drainage-basins:hydrograph}{流量过程线}是一个三角形，底长三天、高为 \qty{1200}{m^3/s}，其
面积为 $\tfrac12 \times 1200 \times 3 \times 86\,400 \approx
\qty{1.56e8}{m^3}$，即 1.56 亿立方米。高于
\qty{800}{m^3/s} 的部分是一个形状相同的较小三角形，其
高度（$400$，而非 $1200$）和历时（一天，而非三天）均缩为原来的三分之一，
面积为 $\tfrac12 \times 400 \times 86\,400 \approx \qty{1.73e7}{m^3}$，
1700 万立方米。这小于水库能够储存的
6000 万立方米，故可以。
\end{solution}

\begin{solution}{exo:g10:drainage-basins:13}
河水满岸时，深度为 $800/(100 \times 2) = \qty{4}{m}$。
（1）河漫滩上\omterm{def:g10:drainage-basins:discharge}{流量}为 \qty{400}{m^3/s}，宽 \qty{900}{m}，流速为
\qty{0.5}{m/s}：其深度为 $400/(900 \times 0.5) \approx \qty{0.9}{m}$。
（2）堤防之间：水深 $1200/(100 \times 2) = \qty{6}{m}$，
高出河岸 \qty{2}{m}，超过无堤防时河漫滩水深的
两倍以上。
\end{solution}

\begin{solution}{exo:g10:drainage-basins:14}
设 $(1 - 1/T)^{30} = 0.5$，于是 $1 - 1/T = 0.5^{1/30} \approx 0.97716$，
$1/T \approx 0.02284$，故 $T \approx 44$ 年。
\end{solution}

\begin{solution}{exo:g10:drainage-basins:15}
每个方格将其水流送往最低的相邻方格（行由北向南，
列由西向东）：
\begin{center}
\begin{tabular}{|c|c|c|c|}
\hline
出口 & 西 & 西 & 西 \\ \hline
北 & 西 & 西 & 南 \\ \hline
北 & 北 & 南 & 南 \\ \hline
北 & 东 & 东 & 出口 \\ \hline
\end{tabular}
\end{center}
沿箭头方向，方格 47、39、44、37、46 和 30 汇入
东南出口：共 6 个方格，\qty{0.06}{km^2}。其余 10 个方格汇入
西北出口：\qty{0.10}{km^2}。\omterm{def:g10:drainage-basins:drainage-basin}{分水岭}从东侧边缘
出发，依次在 50 与 47 之间、48 与 47 之间、48 与 44 之间、
49 与 44 之间、49 与 46 之间，以及 51 与 46 之间，直至南侧边缘。
\end{solution}

\begin{solution}{pb:g10:drainage-basins:1}
\textbf{1.} $11\,560 \times 1.609^2 \approx \qty{29900}{km^2}$.

\textbf{2.} 测站上游\omterm{def:g10:drainage-basins:drainage-basin}{流域}的\omterm{def:g10:drainage-basins:drainage-basin}{分水岭}。

\textbf{3.} $484\,000 \times 0.3048^3 \approx \qty{13700}{m^3/s}$.

\textbf{4.} 约为 $13\,705/3341 \approx 4.1$ 倍。

\textbf{5.} $A = 400 \times 8 = \qty{3200}{m^2}$ 且 $v = Q/A = 13\,705/3200
\approx \qty{4.3}{m/s}$.

\textbf{6.} 约 \qty{330}{m^3/s}。

\textbf{7.} \qty{9231}{m^3/s}，出现于 1 月 21 日。

\textbf{8.} 三天。

\textbf{9.} $9231 \times 86\,400/(29\,900 \times 10^6) \approx
\qty{0.027}{m}$：约 \qty{27}{mm}。

\textbf{10.} 新的水进入河流：\omterm{def:g10:drainage-basins:drainage-basin}{流域}内又一次降雨
（或融雪）带来了第二次较小的径流脉冲。

\textbf{11.} 峰值后十七天。涨水段的径流来自
河流附近坡面的径流，迅速且集中地抵达河道；
退水段时，河流仍接收\omterm{def:g10:drainage-basins:drainage-basin}{流域}最远处的水，
以及缓慢经土壤排出的水。

\textbf{12.} $T = n + 1 = 96$ 年。

\textbf{13.} $96/2 = 48$ 年。

\textbf{14.} $96/3 = 32$ 年。

\textbf{15.} $96/4 = 24$ 年和 $96/5 = 19.2$ 年。两者差别不大：无论哪种算法，
该\omterm{def:g10:drainage-basins:discharge}{流量}大约每二十年重现一次。

\textbf{16.} 五次（1936、1937、1943、1972 和 1996）：$5/95 \approx
0.053$，即每年约 \qty{5}{\percent}，\omterm{def:g10:drainage-basins:return-period}{重现期}约为
$96/5 \approx 19$ 年。

\textbf{17.} $1/100 = 0.01$, \qty{1}{\percent}.

\textbf{18.} $1 - 0.01 = 0.99$.

\textbf{19.} $0.99^{30} \approx 0.740$.

\textbf{20.} $1 - 0.740 = 0.26$：\textbf{30 年内至少发生一次百年一遇洪水的概率约为
\qty{26}{\percent}}，即四分之一。
\end{solution}
